\documentclass[10pt,a4paper]{article}

\usepackage[utf8]{inputenc}
\usepackage[T1]{fontenc}
\usepackage[spanish]{babel}
\usepackage{geometry}
\usepackage{booktabs}
\usepackage{array}
\usepackage{enumitem}
\usepackage{hyperref}
\usepackage{url}

\geometry{
  left=1.55cm,
  right=1.55cm,
  top=1.55cm,
  bottom=1.65cm
}

\hypersetup{
  colorlinks=true,
  linkcolor=black,
  urlcolor=blue
}

\setlength{\parindent}{0pt}
\setlength{\parskip}{0.38em}
\renewcommand{\arraystretch}{1.08}
\sloppy

\title{\textbf{Informe técnico -- Trabajo teórico-práctico}\\
\large Convenio PLUS TI -- Universidad del Valle 2026}
\author{
Andre Marroquín -- Carné: 22266\\
Gabriel Paz -- Carné: 221087\\
Mauricio Lemus -- Carné: 22461
}
\date{}

\begin{document}

\maketitle

\section{Resumen}

El trabajo se desarrolló en dos fases. En la primera fase se trabajó con el dataset
\texttt{01\_bo\_vip\_seed22\_n100000.csv}, correspondiente a transacciones sintéticas de tarjetas
de crédito de un banco de Bolivia con clientela VIP. El objetivo principal fue evaluar modelos 
de clasificación binaria para detectar fraude, dando especial importancia a la reducción de 
falsos positivos entre las alertas generadas. La métrica principal usada para comparar estrategias fue:

\[
\text{ratio de falsos positivos} = \frac{FP}{TP + FP}
\]

Esta métrica mide qué proporción de las alertas marcadas como fraude son realmente falsas alarmas.
Esto es importante en fraude porque un modelo puede tener buen recall y parecer adecuado
en las métricas generales, pero resultar poco útil en la operación si genera demasiadas
alertas incorrectas.

En la segunda fase se trabajó con tres bancos. Banco~1 y Banco~2 tenían etiqueta real
\texttt{is\_fraud} y se utilizaron para entrenamiento, validación y comparación de estrategias.
Banco~3 era el conjunto objetivo que había que predecir, por lo que se
trató como no etiquetado y se utilizó para generar las inferencias. A partir de
Banco~1 y Banco~2 se entrenaron modelos locales por banco, se construyeron ensambles federados
simulados y se compararon contra modelos centralizados.

Los notebooks anexos son:

\begin{itemize}[leftmargin=*]
  \item \texttt{ObjetivoA.ipynb}: implementación completa de métricas custom con LightGBM.
  \item \texttt{ObjetivoB.ipynb}: entrenamiento federado simulado, comparación de estrategias
        e inferencia sobre Banco~3.
\end{itemize}

\section{Metodología}

\subsection{Fase A: métricas custom para reducir falsos positivos}

El dataset de Banco~1 contiene 100{,}003 transacciones, con 4{,}919 fraudes y 95{,}084 operaciones
legítimas. La tasa global de fraude es de 4.92\%, por lo que se trata de un problema desbalanceado.
Se construyó una variable proxy para identificar transacciones de app móvil o canal digital
equivalente, derivada a partir de señales como canal e-commerce, terminal virtual, modo de
entrada POS y condición de comercio electrónico.

El segmento mobile/app quedó compuesto por 46{,}312 transacciones, con una tasa de fraude
de 8.02\%. El resto del dataset tuvo una tasa de 2.24\%. Esta diferencia justificó evaluar
el desempeño global y también el comportamiento específico del modelo en el segmento digital.

La separación se manejó de la siguiente manera: entrenamiento con enero--mayo de 2025 y evaluación con
junio de 2025. El conjunto de entrenamiento quedó con 83{,}817 filas y el de prueba con 16{,}186 filas.

La ingeniería de variables incluyó indicadores temporales, comportamiento por cliente,
contadores recientes y desviaciones de monto:

\begin{itemize}[leftmargin=*, itemsep=0pt]
  \item \texttt{time\_since\_last\_txn\_min}, \texttt{txn\_count\_last\_1h},
        \texttt{txn\_count\_last\_24h}
  \item \texttt{amount\_zscore\_customer}, \texttt{amount\_ratio\_customer\_mean},
        \texttt{amount\_ratio\_customer\_median}
  \item Contadores históricos por cliente, tarjeta y comercio
  \item Variables de canal digital, horario nocturno, fin de semana y transacción internacional
\end{itemize}

Se entrenó un modelo base LightGBM con métricas tradicionales y tres variantes con funciones
\texttt{feval} personalizadas:

\begin{center}
\small
\begin{tabular}{p{0.39\textwidth} p{0.53\textwidth}}
\toprule
\textbf{Métrica custom} & \textbf{Intención} \\
\midrule
\texttt{mobile\_app\_false\_positive\_ratio}
  & Minimizar directamente la proporción de falsos positivos en el segmento mobile/app. \\[4pt]
\texttt{mobile\_app\_recall\_}\newline\texttt{constrained\_fp\_ratio}
  & Reducir falsos positivos penalizando si el recall cae por debajo de 90\%. \\[4pt]
\texttt{mobile\_app\_alert\_quality}
  & Combinar recall, proporción de falsos positivos y volumen de alertas para
    buscar alertas más útiles. \\
\bottomrule
\end{tabular}
\end{center}

El criterio principal de comparación fue alcanzar aproximadamente 90\% de detección y,
bajo esa restricción, seleccionar la estrategia con menor \texttt{FP\,/\,(TP\,+\,FP)}.

\subsection{Fase B: entrenamiento federado simulado e inferencia sobre Banco 3}

Para la segunda fase se cargaron tres datasets correspondientes a bancos de distintos países
y segmentos:

\begin{center}
\begin{tabular}{lll}
\toprule
\textbf{Banco} & \textbf{País} & \textbf{Segmento} \\
\midrule
Banco 1 & Bolivia  & VIP \\
Banco 2 & Brasil   & Privado \\
Banco 3 & Guatemala & Estatal \\
\bottomrule
\end{tabular}
\end{center}

Solo Banco~1 y Banco~2 contaban con etiqueta válida. Banco~3 se utilizó para generar
la predicción final de fraude.

La ingeniería de variables fue la misma para los tres bancos e incluyó: transformaciones numéricas
(\texttt{log\_amount}, ratio respecto al baseline del cliente), indicadores temporales (hora,
día, mes, transacción nocturna, horario hábil), proxies de canal digital (e-commerce, entrada
manual, contactless), y contadores/acumuladores por cliente, tarjeta y comercio (conteo de
transacciones, monto total, monto promedio, ratio monto/promedio, conteo de comercios únicos).

Para controlar la fuga de información se excluyeron identificadores directos de banco, cliente,
tarjeta, comercio y terminal; columnas de país, ciudad y moneda local; y columnas con
más del 95\% de nulos o cardinalidad excesiva. Esto lo que hizo fue obligar al modelo a aprender patrones
transaccionales generalizables y no memorizar la información específica de cada institución.

Las estrategias entrenadas y comparadas fueron:

\begin{center}
\small
\begin{tabular}{p{0.46\textwidth} p{0.46\textwidth}}
\toprule
\textbf{Estrategia} & \textbf{Descripción} \\
\midrule
Modelo local Banco~1
  & ExtraTrees entrenado exclusivamente con datos de Banco~1. \\[3pt]
Modelo local Banco~2
  & ExtraTrees entrenado exclusivamente con datos de Banco~2. \\[3pt]
Ensamble federado (pesos iguales)
  & Promedio simple de probabilidades de los modelos locales. \\[3pt]
Ensamble federado (pesos por desempeño)
  & Promedio ponderado por F1 en validación de cada modelo local. \\[3pt]
Centralizado baseline (ExtraTrees)
  & Modelo entrenado con datos combinados de Banco~1 y Banco~2. \\[3pt]
Centralizado logístico
  & Regresión logística sobre datos combinados. \\[3pt]
Centralizado boosting
  & HistGradientBoosting sobre datos combinados. \\[3pt]
Boosting con hiperparámetros tuneados
  & HistGradientBoosting optimizado con GridSearchCV y validación cruzada. \\
\bottomrule
\end{tabular}
\end{center}

La búsqueda de hiperparámetros exploró combinaciones de \texttt{max\_iter},
\texttt{learning\_rate}, \texttt{max\_leaf\_nodes} y \texttt{l2\_regularization}
mediante \texttt{GridSearchCV} con validación cruzada estratificada.
El mejor estimador se reentrenó con todos los datos disponibles antes de competir
en la selección final.

La selección de la metodología final se realizó mediante puntuaciones compuestas que combinaban
F1 promedio entre bancos, F1 mínimo entre bancos, ROC AUC, recall y penalización por
false positive ratio. La estrategia que mejor resultado dió en la puntuación se usó para la
inferencia final sobre el 100\% de Banco~3.

\section{Descripción de la implementación práctica}

La Fase~A está hecha en \texttt{ObjetivoA.ipynb}. El notebook normaliza nombres
de columnas, valida el target, detecta columnas relevantes del estándar ISO~8583,
construye la fecha transaccional desde \texttt{de7\_transmission\_datetime}, separa
enero--mayo contra junio, crea variables predictivas y entrena LightGBM. La evaluación
se hizo con umbral fijo 0.50 y con un umbral elegido automáticamente para alcanzar
recall cercano a 90\%, que fue el punto principal de comparación.

La Fase~B está hecha en \texttt{ObjetivoB.ipynb}. El notebook carga los tres
bancos, genera variables equivalentes,
excluye columnas con riesgo de fuga y entrena las estrategias locales, federadas
simuladas y centralizadas. Los modelos se guardaron en
\nolinkurl{outputs_fase_B/04_models}, los reportes en
\nolinkurl{outputs_fase_B/07_reports} y las inferencias en
\nolinkurl{outputs_fase_B/05_predictions}.

\section{Análisis de resultados}

\subsection{Fase A}

El modelo seleccionado como final fue el baseline LightGBM con métricas tradicionales.
Al buscar el umbral que garantiza recall cercano a 90\% en el segmento mobile/app,
el baseline presentó el menor ratio de falsos positivos de las estrategias evaluadas.

\begin{center}
\scriptsize
\begin{tabular}{p{0.36\textwidth}rrrrr}
\toprule
\textbf{Estrategia} & \textbf{Recall} & \textbf{Prec.} & \textbf{FP} & \textbf{TP} & \textbf{Ratio FP} \\
\midrule
Baseline LightGBM & 95.31\% & 14.34\% & 3{,}034 & 508 & 85.66\% \\
\texttt{mobile\_app\_alert\_quality} & 95.31\% & 12.81\% & 3{,}458 & 508 & 87.19\% \\
\texttt{mobile\_app\_recall\_constrained\_fp\_ratio} & 95.50\% & 12.68\% & 3{,}506 & 509 & 87.32\% \\
Modelo tuneado custom & 95.50\% & 11.78\% & 3{,}813 & 509 & 88.22\% \\
\texttt{mobile\_app\_false\_positive\_ratio} & 100.00\% & 7.21\% & 6{,}864 & 533 & 92.79\% \\
\bottomrule
\end{tabular}
\end{center}

Las métricas custom mejoraron el ordenamiento del modelo. Por ejemplo,
\texttt{mobile\_app\_recall\_constrained\_fp\_ratio} alcanzó AUC PR de 0.885 en
mobile/app frente a 0.771 del baseline. Sin embargo, al forzar una detección cerca
a 90\%, ese mejor ranking no se convirtió en menor proporción de falsos positivos.
Con umbral fijo 0.50 sí aparecieron alertas más limpias. La métrica
\texttt{mobile\_app\_recall\_constrained\_fp\_ratio} logró 82.36\% de recall,
93.40\% de precisión y 6.60\% de ratio FP en mobile/app, pero no alcanzó el recall
establecido anteriormente para la comparación principal.

\subsection{Fase B}

Los modelos locales de Banco~1 y Banco~2 se usaron como base para los ensambles federados
simulados. Para la selección final se compararon seis estrategias candidatas sobre los
conjuntos de validación de Banco~1 y Banco~2. La estrategia final se eligió con una
puntuación compuesta que pondera F1 promedio, F1 mínimo entre bancos, ROC AUC, recall
y penaliza el ratio de falsos positivos.

\begin{center}
\scriptsize
\begin{tabular}{p{0.35\textwidth}rrrrrr}
\toprule
\textbf{Estrategia} & \textbf{Umbral} & \textbf{F1 prom.} & \textbf{F1 mín.} & \textbf{ROC AUC} & \textbf{Recall} & \textbf{Ratio FP} \\
\midrule
HistGradientBoosting central & 0.85 & 0.801 & 0.796 & 0.878 & 0.687 & 0.041 \\
HistGradientBoosting tuneado & 0.85 & 0.794 & 0.784 & 0.880 & 0.674 & 0.033 \\
ExtraTrees centralizado & 0.80 & 0.676 & 0.623 & 0.872 & 0.603 & 0.223 \\
Regresión logística central & 0.90 & 0.592 & 0.572 & 0.867 & 0.456 & 0.157 \\
Federado ponderado & 0.65 & 0.556 & 0.409 & 0.864 & 0.523 & 0.406 \\
Federado pesos iguales & 0.60 & 0.544 & 0.393 & 0.863 & 0.556 & 0.467 \\
\bottomrule
\end{tabular}
\end{center}

La mejor estrategia desarrollada fue \texttt{central\_hist\_gradient\_boosting\_model}
con umbral 0.85. Los ensambles federados simulados permitieron comparar transferencia
entre bancos, pero no superaron al mejor modelo centralizado. El federado ponderado
mejoró al federado de pesos iguales en ratio FP, aunque mantuvo menor estabilidad entre
bancos.

Las inferencias sobre el primer 30\% de Banco~3 se generaron con tres
estrategias (baseline centralizado, ensamble con pesos iguales y ensamble ponderado)
y quedan en \nolinkurl{outputs_fase_B/05_predictions}. Para la inferencia final sobre
el 100\% de Banco~3 se usó la estrategia ganadora, con 1{,}739 fraudes predichos sobre
100{,}000 transacciones, equivalente a una tasa esperada de 1.74\%.

\section{Dataset 3 con inferencia realizada}

La inferencia final para Banco~3 quedó guardada en:

\begin{center}
\nolinkurl{outputs_fase_B/05_predictions/bank_3_full_final_inference.xlsx}
\end{center}

El archivo contiene una única columna llamada \texttt{is\_fraud}, con valores
\texttt{True} o \texttt{False}, y mantiene las 100{,}000 transacciones en el orden de
entrada. También se generó el archivo de auditoría:

\begin{center}
\nolinkurl{outputs_fase_B/05_predictions/bank_3_full_final_inference_audit.csv}
\end{center}

Ese archivo de auditoría no es el final, pero permite revisar la
probabilidad asignada, el umbral utilizado y la estrategia que generó cada predicción.

\section{Conclusiones}

El trabajo permitió ver dos enfoques para la detección de fraude transaccional.
En la Fase~A se evidenció que el uso de funciones \texttt{feval} personalizadas en LightGBM
permite ajustar el comportamiento del modelo hacia criterios operativos específicos,
aunque su efecto depende del régimen de umbral bajo el cual se evalúe el modelo.
La división temporal del conjunto de prueba permitió realizar una evaluación más realista y 
evitar fuga de información.

En la Fase~B se trabajó con una simulación de aprendizaje federado para representar un
escenario con varios bancos, sin necesidad de reunir todos los datos originales en un solo
lugar. Los ensambles federados ayudaron a comparar distintas estrategias y a no depender
únicamente de la información de un banco. Sin embargo, con los resultados obtenidos, 
el modelo centralizado HistGradientBoosting terminó siendo la mejor opción para 
generar la inferencia final.

Las principales limitaciones son: (1)~el carácter sintético de los datasets, que puede
no reflejar con exactitud la distribución de fraude real; (2)~la simulación de federación
en un único entorno local, sin comunicación real de parámetros entre nodos; y
(3)~la ausencia de etiqueta confiable en Banco~3, que impide evaluar la calidad real
de las inferencias generadas.

\section{Anexos del repositorio}

\begin{itemize}[leftmargin=*, itemsep=0pt]
  \item Notebook Fase A: \texttt{ObjetivoA.ipynb}
  \item Notebook Fase B: \texttt{ObjetivoB.ipynb}
  \item Resultados Fase A: \nolinkurl{outputs_fase_A/07_final_model/final_metrics.csv}
  \item Resumen Fase A: \nolinkurl{outputs_fase_A/07_final_model/final_model_summary.json}
  \item Resultados Fase B: \nolinkurl{outputs_fase_B/07_reports/selection_table.csv}
  \item Resumen Fase B: \nolinkurl{outputs_fase_B/07_reports/expected_result_summary.json}
  \item Inferencia final Banco~3: \nolinkurl{outputs_fase_B/05_predictions/bank_3_full_final_inference.xlsx}
\end{itemize}

\end{document}
